\documentclass[dvipdfmx,ja=standard]{beamer}

\usetheme[block=fill]{metropolis} % テーマの指定

% === 文字化け・文字消失対策のための追加設定 ===
\usepackage[deluxe,uplatex]{otf} % 高機能な日本語フォント処理
\renewcommand{\kanjifamilydefault}{\gtdefault} % 日本語の既定をゴシック体に
\font\zero=cmr10 at 0pt % Metropolisのフォントバグを回避するおまじない
% =============================================

\definecolor{katsuobushi}{HTML}{cd853f}

\title{テストスライド}
\author{俺}
\date{\today}

\begin{document}

% 最初のタイトルスライド
\begin{frame}
    \maketitle
\end{frame}

\begin{frame}
    \setbeamercolor{block title}{bg=katsuobushi}
    \setbeamercolor{block body}{bg=katsuobushi}
    \setbeamercolor{block title alerted}{bg=katsuobushi}
    \setbeamercolor{block body alerted}{bg=katsuobushi}
    \setbeamercolor{block title example}{bg=katsuobushi}
    \setbeamercolor{block body example}{bg=katsuobushi}
    
    \frametitle{一枚目のスライド}
    普通に文中で\LaTeX コマンドが使用できます。（そらそう）
    \begin{block}{ブロックのタイトル}
        block環境を作成しbeginの第二引数にタイトルを渡すとこうなります。
    \end{block}
    \begin{alertblock}{アラート}
        alertblockを指定するとこんな風になります。
    \end{alertblock}
    \begin{exampleblock}{例}
        exampleblockを指定するとこうなります。
    \end{exampleblock}
\end{frame}

\begin{frame}
    \frametitle{横並べ}
    スライドで左右比較するときなどにこのようにします
    \begin{columns}
        \begin{column}{.4\linewidth}
            \begin{block}{内容1}
                内容1内容1内容1内容1内容1内容1内容1内容1内容1内容1内容1内容1内容1内容1
            \end{block}
        \end{column}
        \begin{column}{.4\linewidth}
            \begin{block}{内容2}
                内容2内容2内容2内容2内容2内容2内容2内容2内容2内容2内容2内容2内容2内容2
            \end{block}
        \end{column}
    \end{columns}
    このとき各columnの幅は合計100\%にしてしまうとデザイン的にあんまよくないです
\end{frame}

\end{document}
